% introduction.tex

\chapter{Introduction}

\section{A Shift in Perspective}

Railway alignment is commonly described as a geometric curve
\(\gamma(s)\) embedded in space.

This view is sufficient for visualization, but it fails to capture the
quantities that actually govern railway design:
curvature, cant, dynamics, and realization.

In practice, railway tracks are not designed, constructed, or maintained
as curves, but through operations acting on curvature and discrete
control points.

\begin{summary}
Alignment is not a curve.
It is a curvature-driven system whose geometry emerges from its
generative structure.
\end{summary}

\section{Limitations of Current Practice}

Despite a long tradition of railway engineering, alignment modelling
remains fragmented across several domains:

\begin{itemize}
\item geometric construction of curves,
\item rule-based engineering design,
\item data exchange formats such as LandXML and IFC,
\item numerical evaluation and simulation.
\end{itemize}

These aspects are typically handled in separate tools and workflows,
leading to discontinuities between modelling, analysis, and design.

In particular:
\begin{itemize}
\item geometry is treated independently of dynamics,
\item transition curves are selected from predefined catalogues,
\item data formats focus on representation rather than computation,
\item optimization is not part of the core modelling paradigm.
\end{itemize}

As a consequence, alignment design lacks a unified internal model that
connects data, geometry, dynamics, and realization.

\section{Problem Statement}

The absence of a unified modelling framework leads to fundamental
limitations:

\begin{itemize}
\item inconsistent interpretation of alignment data,
\item separation of geometry and dynamics,
\item limited support for optimization-based design,
\item high effort for integrating heterogeneous data sources.
\end{itemize}

These issues become increasingly critical as infrastructure systems grow
in complexity and demand higher levels of automation, precision, and
traceability.

\section{AIM: A Unified Framework}

This work introduces \textbf{Alignment-Based Information Modelling (AIM)}
as a framework addressing these limitations.

The central idea is to represent alignment as a structured, state-based
object defined by curvature and cant:
\[
(\kappa(s), \phi(s)),
\]
from which all geometric and dynamic quantities are derived.

Within this framework:
\begin{itemize}
\item curvature is the primary modelling variable,
\item geometry is obtained by integration,
\item transition curves are interpreted as control functions,
\item alignment becomes a generative system rather than a static object.
\end{itemize}

\section{Core Claim}

\begin{proposition}
Railway alignment modelling can be reformulated as a curvature-driven,
realization-aware, and optimization-ready system.
\end{proposition}

This reformulation enables a unified treatment of:
\begin{itemize}
\item geometry,
\item dynamics,
\item construction and realization,
\item data integration,
\item and optimization.
\end{itemize}

\section{Consequences}

Adopting this perspective leads to several key consequences:

\begin{itemize}
\item alignment design becomes an optimization problem,
\item realization is treated as an operator acting on the model,
\item data reconstruction is formulated as an inverse problem,
\item hybrid models bridge parametric structure and measured data,
\item analytical derivatives arise naturally from the formulation.
\end{itemize}

\section{Scope of the Work}

The present manuscript develops the AIM framework across the following
dimensions:

\begin{itemize}
\item mathematical foundations of curvature-based modelling,
\item structured representation of alignment elements,
\item classification and interpretation of transition curves,
\item modelling of realization and construction processes,
\item formulation of optimization problems and solution methods,
\item integration of real-world data and coordinate systems,
\item application to railway engineering use cases.
\end{itemize}

\section{Contribution}

The main contributions of this work are:

\begin{itemize}
\item a unified state-based alignment model,
\item the reinterpretation of transition curves as curvature control functions,
\item the introduction of a realization-aware modelling perspective,
\item the formulation of alignment design as a structured optimization problem,
\item the development of a pipeline connecting data, model, and computation.
\end{itemize}

Together, these contributions establish a coherent framework linking
theoretical modelling and practical engineering.

\section{Structure of the Manuscript}

The manuscript is organized into several parts, progressing from
conceptual foundations to practical applications:

\begin{itemize}
\item motivation and conceptual shift,
\item mathematical and state-based foundations,
\item modelling and representation,
\item realization and system perspective,
\item optimization methods,
\item engineering constraints and standards,
\item applications and outlook.
\end{itemize}

\section{Outlook}

The AIM framework is intended as a foundation for future development.

It opens the path toward:
\begin{itemize}
\item optimization-driven design systems,
\item deeper integration with BIM,
\item network-level alignment modelling,
\item interactive and real-time engineering tools.
\end{itemize}

In this sense, the work contributes to a broader transition from static
geometry toward computational infrastructure modelling.
